\section*{Capítulo \ref{ch:g7:atoms-and-molecules} --- Átomos e moléculas}

\begin{solution}{exo:g7:atoms-and-molecules:1}
Um \omterm{def:g7:atoms-and-molecules:atom}{átomo} é uma das partículas extremamente pequenas a partir das quais
toda a matéria é construída (existem cerca de cem tipos). Uma \omterm{def:g7:atoms-and-molecules:molecule}{molécula} é
um grupo de \omterm{def:g7:atoms-and-molecules:atom}{átomos} firmemente unidos, o menor pedaço de uma substância
que ainda é essa substância.
\end{solution}

\begin{solution}{exo:g7:atoms-and-molecules:2}
\ce{H}, \ce{C}, \ce{N}, \ce{O}, \ce{Na}, \ce{Fe}.
\end{solution}

\begin{solution}{exo:g7:atoms-and-molecules:3}
1 \omterm{def:g7:atoms-and-molecules:atom}{átomo} de carbono e 2 \omterm{def:g7:atoms-and-molecules:atom}{átomos} de oxigênio: 3 \omterm{def:g7:atoms-and-molecules:atom}{átomos} ao todo.
\end{solution}

\begin{solution}{exo:g7:atoms-and-molecules:4}
Água, oxigênio, metano, nitrogênio.
\end{solution}

\begin{solution}{exo:g7:atoms-and-molecules:5}
\ce{NO2}.
\end{solution}

\begin{solution}{exo:g7:atoms-and-molecules:6}
\ce{O2} é uma \omterm{def:g7:atoms-and-molecules:molecule}{molécula} formada por dois \omterm{def:g7:atoms-and-molecules:atom}{átomos} de oxigênio ligados;
\ce{2O} são dois \omterm{def:g7:atoms-and-molecules:atom}{átomos} de oxigênio separados. \ce{CO} é uma \omterm{def:g7:atoms-and-molecules:molecule}{molécula} de
um \omterm{def:g7:atoms-and-molecules:atom}{átomo} de carbono e um de oxigênio (duas letras maiúsculas, dois
\omterm{def:g7:atoms-and-molecules:symbol}{símbolos}); \ce{Co} é um \omterm{def:g7:atoms-and-molecules:atom}{átomo} de cobalto (um \omterm{def:g7:atoms-and-molecules:symbol}{símbolo}, maiúscula seguida
de minúscula).
\end{solution}

\begin{solution}{exo:g7:atoms-and-molecules:7}
\ce{3H2O}: $3\times 2 = 6$ \omterm{def:g7:atoms-and-molecules:atom}{átomos} de hidrogênio e $3\times 1 = 3$ \omterm{def:g7:atoms-and-molecules:atom}{átomos}
de oxigênio. \ce{5O2}: nenhum hidrogênio, $5\times 2 = 10$ \omterm{def:g7:atoms-and-molecules:atom}{átomos} de oxigênio.
\end{solution}

\begin{solution}{exo:g7:atoms-and-molecules:8}
Preto é carbono, vermelho é oxigênio: um carbono, dois oxigênios,
\ce{CO2}, dióxido de carbono.
\end{solution}

\begin{solution}{exo:g7:atoms-and-molecules:9}
A água de Dalton era \ce{HO}. A \omterm{def:g7:atoms-and-molecules:molecule}{molécula} real, \ce{H2O}, tem um \omterm{def:g7:atoms-and-molecules:atom}{átomo} de
hidrogênio a mais: faltava um \omterm{def:g7:atoms-and-molecules:atom}{átomo} na \omterm{def:g7:atoms-and-molecules:molecule}{molécula} de Dalton.
\end{solution}

\begin{solution}{exo:g7:atoms-and-molecules:10}
A água pura não é uma \omterm{def:g6:pure-substances-and-mixtures:mixture}{mistura}: todas as suas \omterm{def:g7:atoms-and-molecules:molecule}{moléculas} são \omterm{def:g7:atoms-and-molecules:molecule}{moléculas} de
\ce{H2O} idênticas. O \omterm{def:g4:air-a-mixture-of-gases:air}{ar} é uma \omterm{def:g6:pure-substances-and-mixtures:mixture}{mistura}: ele contém \omterm{def:g7:atoms-and-molecules:molecule}{moléculas} de
nitrogênio, \omterm{def:g7:atoms-and-molecules:molecule}{moléculas} de oxigênio, \omterm{def:g7:atoms-and-molecules:atom}{átomos} de argônio e algumas \omterm{def:g7:atoms-and-molecules:molecule}{moléculas}
de dióxido de carbono.
\end{solution}

\begin{solution}{exo:g7:atoms-and-molecules:11}
Dez milhões de \omterm{def:g7:atoms-and-molecules:atom}{átomos} formam \qty{1}{mm}, então \qty{1}{cm} $= \qty{10}{mm}$
exige $10 \times$ dez milhões $=$ cem milhões de \omterm{def:g7:atoms-and-molecules:atom}{átomos}.
$\qty{100}{m} = \num{100000}$~mm exige $\num{100000}\times$ dez milhões
$=$ um milhão de milhões de \omterm{def:g7:atoms-and-molecules:atom}{átomos}.
\end{solution}

\begin{solution}{exo:g7:atoms-and-molecules:12}
$6 + 12 + 6 = 24$ \omterm{def:g7:atoms-and-molecules:atom}{átomos}. Hidrogênio: $\frac{12}{24} = \qty{50}{\percent}$.
Carbono: $\frac{6}{24} = \qty{25}{\percent}$ (e o oxigênio, os outros
\qty{25}{\percent}).
\end{solution}

\begin{solution}{pb:g7:atoms-and-molecules:1}
\textbf{1.} Nitrogênio \ce{N2}, 2 \omterm{def:g7:atoms-and-molecules:atom}{átomos}; oxigênio \ce{O2}, 2 \omterm{def:g7:atoms-and-molecules:atom}{átomos};
dióxido de carbono \ce{CO2}, 3 \omterm{def:g7:atoms-and-molecules:atom}{átomos}.

\textbf{2.} Os \omterm{def:g7:atoms-and-molecules:atom}{átomos} de argônio não se juntam em \omterm{def:g7:atoms-and-molecules:molecule}{moléculas}: o argônio
é formado por \omterm{def:g7:atoms-and-molecules:atom}{átomos} isolados.

\textbf{3.} Nitrogênio: $78 \times 2 = 156$ \omterm{def:g7:atoms-and-molecules:atom}{átomos}. Oxigênio: $21\times 2 =
42$ \omterm{def:g7:atoms-and-molecules:atom}{átomos}.

\textbf{4.} $156 + 42 + 1 = 199$ \omterm{def:g7:atoms-and-molecules:atom}{átomos}.

\textbf{5.} $\frac{156}{199} \approx 0.78$, cerca de \qty{78}{\percent}.

\textbf{6.} Nas \omterm{def:g7:atoms-and-molecules:molecule}{moléculas} de oxigênio: $16\times 2 = 32$; nas \omterm{def:g7:atoms-and-molecules:molecule}{moléculas}
de dióxido de carbono: $4 \times 2 = 8$; ao todo $32 + 8 = 40$.

\textbf{7.} Faltam $42 - 40 = 2$ \omterm{def:g7:atoms-and-molecules:atom}{átomos} de oxigênio: o equivalente a uma
\omterm{def:g7:atoms-and-molecules:molecule}{molécula} de oxigênio. Eles foram para \omterm{def:g7:atoms-and-molecules:molecule}{moléculas} de água, \ce{H2O}, que o
corpo fabrica com o hidrogênio dos alimentos; essa água sai em parte
como vapor de água no \omterm{def:g4:air-a-mixture-of-gases:air}{ar} expirado (retirado destes números).

\textbf{8.} No \omterm{def:g4:air-a-mixture-of-gases:air}{ar} inspirado: 0. No \omterm{def:g4:air-a-mixture-of-gases:air}{ar} expirado: 4 (um em cada \ce{CO2}).

\textbf{9.} Dos alimentos que o corpo usa: os açúcares e as gorduras
contêm \omterm{def:g7:atoms-and-molecules:atom}{átomos} de carbono, e o corpo os libera no dióxido de carbono que
expira.

\textbf{10.} Desapareceram $21 - 16 = 5$ \omterm{def:g7:atoms-and-molecules:molecule}{moléculas} de oxigênio e
apareceram 4 \omterm{def:g7:atoms-and-molecules:molecule}{moléculas} de dióxido de carbono.

\textbf{11.} Nitrogênio: duas bolas azuis ligadas por três varetas.
Oxigênio: duas bolas vermelhas ligadas por duas varetas. Dióxido de
carbono: uma bola preta entre duas bolas vermelhas, cada uma ligada a ela
por duas varetas, em linha reta. Argônio: uma única bola azul-clara.

\textbf{12.} $156 + 32 + 12 + 1 = 201$ \omterm{def:g7:atoms-and-molecules:atom}{átomos} (as 4 \omterm{def:g7:atoms-and-molecules:molecule}{moléculas} de dióxido
de carbono contêm 12 \omterm{def:g7:atoms-and-molecules:atom}{átomos}). São $201 - 199 = 2$ \omterm{def:g7:atoms-and-molecules:atom}{átomos} a mais do que no
\omterm{def:g4:air-a-mixture-of-gases:air}{ar} inspirado: 4 \omterm{def:g7:atoms-and-molecules:atom}{átomos} de carbono ganhos dos alimentos, menos os 2 \omterm{def:g7:atoms-and-molecules:atom}{átomos}
de oxigênio que saíram na água.
\end{solution}
